\section{The infimum: quasi-isomorphic complexes and moduli}\label{sec:moduli}

The charges are ordered. This section identifies the infimum of the phase
potential and shows that, up to an explicit constant, it is the same for
quasi-isomorphic complexes of the class \eqref{eq:complex}.

\subsection{Volume data}
A \emph{volume datum} for a complex $\mathsf E$ is a triple
$\varpi=(\varpi_V,\varpi_U,\varpi_W)$ of nonzero elements of the top exterior
powers $\det V,\det U,\det W$. If these are the wedges of the reference bases,
then by \eqref{eq:primitive}
\begin{equation}\label{eq:Kvol}
 \cK(h;\varpi)=\theta_0\bigl(\log\norm{\varpi_V}^2_{h_V}-\log\norm{\varpi_U}^2_{h_U}\bigr)
 -\theta_1\log\norm{\varpi_W}^2_{h_W}+\theta_0\tr\Ein(B)+\tr\bigl(CR(B)\bigr),
\end{equation}
and we use \eqref{eq:Kvol} as the definition of $\cK$ for an arbitrary volume
datum. Let $H=\ker D$. There is a canonical isomorphism
\[
 \iota_D:\det V\otimes(\det U)^{-1}\to\det H,\qquad
 (v_1\wedge\dots\wedge v_{n_0+r})\otimes(Dv_{n_0+1}\wedge\dots\wedge Dv_{n_0+r})^{-1}
 \mapsto v_1\wedge\dots\wedge v_{n_0},
\]
for any basis $(v_i)$ of $V$ whose first $n_0$ vectors span $H$. Put
$\bar\varpi=(\iota_D(\varpi_V\otimes\varpi_U^{-1}),\varpi_W)$, a volume datum
for the cohomology representation $\mathcal H=(H\to W)$, and
\begin{gather*}
 \cK_{\coh}(\mathcal H;h_0,h_1;\bar\varpi)=q+\theta_0\log\norm{\iota_D(\varpi_V\otimes\varpi_U^{-1})}^2_{h_0}
 -\theta_1\log\norm{\varpi_W}^2_{h_1},\\
 \cK_*(\mathcal H;\bar\varpi)=\inf_{h_0,h_1}\cK_{\coh}.
\end{gather*}

\begin{lemma}\label{lem:determinant}
For all metrics, with $h_0=h_V|_H$ and $\det{}'$ the product of the nonzero
eigenvalues,
\[
 \log\norm{\varpi_V}^2_{h_V}-\log\norm{\varpi_U}^2_{h_U}+\log\det{}'(D^\dagger D)
 =\log\norm{\iota_D(\varpi_V\otimes\varpi_U^{-1})}^2_{h_0}.
\]
\end{lemma}

\begin{proof}
Take an $h_V$-orthonormal basis $(v_i)$ with $v_1,\dots,v_{n_0}\in H$ and an
$h_U$-orthonormal basis $(u_i)$, let $F$ be the matrix of $D|_{H^\perp}$, and
write $\varpi_V=a_Vv_1\wedge\dots\wedge v_{n_0+r}$,
$\varpi_U=a_Uu_1\wedge\dots\wedge u_r$. Then
$\iota_D(\varpi_V\otimes\varpi_U^{-1})=(a_V/a_U)\det F\cdot v_1\wedge\dots\wedge v_{n_0}$,
and $|\det F|^2=\det{}'(D^\dagger D)$.
\end{proof}

This is the comparison already used in the proof of
Proposition~\ref{prop:decomposition}, where $c_D$ is the logarithm of the
squared modulus of $\iota_D$ in the reference bases.

\begin{theorem}[Infimum of the potential]\label{thm:infimum}
Let the charges be ordered. For every complex and volume datum,
\begin{gather*}
 \cK(h;\varpi)=\cK_{\coh}(\mathcal H;h_0,h_1;\bar\varpi)
 +\theta_0\bigl[r\gamma+\tr{}'E_1(D^\dagger D)\bigr]+\tr Y,\\
 \inf_h\cK(h;\varpi)=\cK_*(\mathcal H;\bar\varpi)+\theta_0\gamma\,r,
\end{gather*}
where $\tr{}'$ is the trace over the nonzero spectrum; the infimum is not
attained. Consequently, let $\mathsf E,\mathsf E'$ be two complexes of the
class \eqref{eq:complex} with semistable cohomology and
$f=(f_V,f_U,f_W):\mathsf E\to\mathsf E'$ a quasi-isomorphism of complexes of
representations, that is, a morphism inducing an isomorphism
$\mathcal H\to\mathcal H'$ of the cohomology representations. If the volume
data correspond, $\bar\varpi'=f_*\bar\varpi$, then
\[
 \inf_{h'}\cK_{\mathsf E'}(h';\varpi')-\inf_h\cK_{\mathsf E}(h;\varpi)=\theta_0\gamma\,(r'-r).
\]
\end{theorem}

\begin{proof}
The first identity is Proposition~\ref{prop:decomposition} with
Lemma~\ref{lem:determinant} and
\[
 \tr\Ein(D^\dagger D)=\log\det{}'(D^\dagger D)+r\gamma+\tr{}'E_1(D^\dagger D);
\]
the infimum is Corollary~\ref{cor:infimum}. The function $\cK_*$ is
transported by isomorphisms of representations with volume data.
\end{proof}

We note three special cases. For an invertible triple $g=(g_V,g_U,g_W)$ and fixed
$\varpi$,
$\inf\cK_{g\cdot\mathsf E}=\inf\cK_{\mathsf E}-\theta_0\log|\det g_V/\det g_U|^2+\theta_1\log|\det g_W|^2$.
A shear $g_V=\left(\begin{smallmatrix}\id&S\\0&\id\end{smallmatrix}\right)$ in
a splitting $V=H\oplus H'$, which leaves $D$ and the arrows on $H$ unchanged
and replaces the arrows $T_j$ on $H'$ by $T_j-A_j|_HS$, does not change the
infimum. Adding an acyclic summand $H''\xrightarrow{F''}U''$ with arbitrary
arrows $H''\to W$ adds
$\theta_0[\gamma\dim U''+\log|\det F''|^2]$, independently of those arrows.
In \cite[\S8.1]{wue26} the shear and the removal of orthogonal acyclic
summands were compared for fixed metrics on $V,U,W$ and $\lambda$ large.

\begin{remark}\label{rem:normalization}
The \emph{normalized potential} $\cK-\theta_0\gamma\,r$ has infimum
$\cK_*(\mathcal H;\bar\varpi)$, the infimum of the Kempf--Ness function of the
cohomology representation, for every complex with that cohomology. The
constant reflects the choice of regularization: for $b>0$,
$\log b=\int_0^1(1-e^{-xb})\frac{dx}x-\int_1^\infty e^{-xb}\frac{dx}x-\gamma$,
and $\Ein(b)$ is the first integral alone. The primitive of $\Omega$
is determined on the metrics of one complex only up to an additive constant,
and the normalized potential fixes these constants compatibly.
\end{remark}

\subsection{The potential on moduli}
Let $\mathcal R=\operatorname{Hom}(\CC^{n_0},\CC^{n_1})^k$ be the space of
representations of dimension vector $(n_0,n_1)$, with the action of
$\mathrm G=\mathrm{GL}_{n_0}\times\mathrm{GL}_{n_1}$,
$b_j\mapsto g_1b_jg_0^{-1}$, and $\mathrm K=\mathrm U(n_0)\times\mathrm U(n_1)$.
Put
\[
 \varkappa(b)=\inf_{g\in\mathrm G}\Bigl[\norm{g\cdot b}^2+\theta_0\log|\det g_0|^2
 -\theta_1\log|\det g_1|^2\Bigr],
\]
let $\mathcal R^{\mathrm s}$ be the set of stable points,
$\mathcal R_0=\{C_0=\theta_0\id,\ C_1=\theta_1\id\}\cap\mathcal R^{\mathrm s}$
the zeros of the moment map in it, and
$\mathcal M^{\mathrm s}=\mathcal R^{\mathrm s}/\mathrm G=\mathcal R_0/\mathrm K$
the moduli space of stable representations \cite{kin94}, with the K\"ahler
form $\omega_\theta$ obtained by reduction of
$\omega_{\mathrm{flat}}=\ii\partial\bar\partial\norm b^2$: its pullback to
$\mathcal R_0$ is the restriction of $\omega_{\mathrm{flat}}$. Let
$\mathrm{pr}:\mathcal R^{\mathrm s}\to\mathcal M^{\mathrm s}$ be the
projection. The stable locus may be empty, for instance for $k=2$ and
$n_0=n_1\geq2$; the statements about it are then vacuous.

\begin{lemma}\label{lem:kappa}
The function $\varkappa$ has the following properties.
\begin{enumerate}
\item[(a)] $\varkappa(g\cdot b)=\varkappa(b)-\theta_0\log|\det g_0|^2+\theta_1\log|\det g_1|^2$.
\item[(b)] $\varkappa$ is real-analytic on $\mathcal R^{\mathrm s}$ and
 $\ii\partial\bar\partial\varkappa=\mathrm{pr}^*\omega_\theta$ there.
\end{enumerate}
\end{lemma}

\begin{proof}
(a) Substitute $g'g$ in the infimum. (b) Let $b_0\in\mathcal R_0$ and, for a
pair $\zeta=(\zeta_0,\zeta_1)$ of Hermitian matrices, let
$f(\zeta,b)=\norm{e^\zeta\cdot b}^2+2\theta_0\tr\zeta_0-2\theta_1\tr\zeta_1$,
so $\varkappa(b)=\inf_\zeta f(\zeta,b)$. Then $\partial_\zeta f(0,b_0)=0$, and
$\partial_\zeta^2f(0,b_0)[\zeta,\zeta]=4\norm{\zeta\cdot b_0}^2$, with
$\zeta\cdot b=(\zeta_1b_j-b_j\zeta_0)_j$ the infinitesimal action, is positive
definite modulo the scalars $(c,c)$, on which $f$ is constant. By the
implicit function theorem there is a critical point $\zeta(b)$ of
$f(\cdot,b)$, real-analytic near $b_0$. It is the global minimizer:
$b'=e^{\zeta(b)}\cdot b$ lies in $\mathcal R_0$, the function
$x\mapsto f(x\zeta',b')$ is convex with zero derivative at $0$ for every
$\zeta'$, and every element of $\mathrm G$ is $ue^{\zeta'}$ with
$u\in\mathrm K$. For a holomorphic curve $b(w)$ with $b(0)=b_0$ and
$\dot b=b'(0)$, differentiating the minimum gives
$\partial_w\partial_{\bar w}\varkappa
=f_{w\bar w}-f_{w\zeta}f_{\zeta\zeta}^{-1}f_{\zeta\bar w}$ with
$f_{w\bar w}=\norm{\dot b}^2$,
$f_{w\zeta}[\zeta]=2\inner{\zeta\cdot b_0}{\dot b}$ and
$f_{\zeta\zeta}[\zeta,\zeta']=4\Rea\inner{\zeta\cdot b_0}{\zeta'\cdot b_0}$. At
a point of $\mathcal R_0$ the $\mathrm K$-orbit is isotropic, so the real
space $\{\zeta\cdot b_0\}$ is orthogonal to its multiple by $\ii$, and a real
orthonormal basis of it is a unitary basis of the tangent space
$\mathfrak g\cdot b_0$ of the complex orbit. Hence
\[
 \partial_w\partial_{\bar w}\varkappa(b(w))\big|_{w=0}
 =\norm{\dot b}^2-\norm{\Pi_{\mathfrak g\cdot b_0}\dot b}^2,
\]
the squared length of the component of $\dot b$ orthogonal to the complex
orbit, which is the reduced form evaluated on $d\,\mathrm{pr}(\dot b)$. For a
general stable point: every stable orbit meets $\mathcal R_0$
\cite{kn79,kin94}, and for fixed $g$ the translation $b\mapsto g\cdot b$ is a
biholomorphism of $\mathcal R^{\mathrm s}$ commuting with $\mathrm{pr}$ and
changing $\varkappa$ by a constant, by (a); so both sides of the identity are
invariant under it.
\end{proof}

Lemma~\ref{lem:kappa}(b) is the classical description of the K\"ahler
potential of a K\"ahler quotient through the shifted ambient potential; for
tori see \cite{bg04}. We include the short proof for the reader's
convenience.

Let now $\mathcal P$ be a complex manifold with holomorphic maps
$p\mapsto D(p),A_j(p)$ (matrices in fixed bases), $D(p)$ surjective, $\varpi$
the volume datum of the bases, $\mathsf E(p)$ the corresponding complexes,
$\cK_\infty(p)=\inf_h\cK_{\mathsf E(p)}(h;\varpi)$, and $\mathcal P^{\mathrm s}$
the open set where the cohomology representation is stable.

\begin{theorem}[The potential on moduli]\label{thm:moduli}
Let the charges be ordered. On $\mathcal P^{\mathrm s}$ the function
$\cK_\infty$ is real-analytic and
\[
 \ii\partial\bar\partial\cK_\infty=\mathrm{coh}^*\omega_\theta,
\]
where $\mathrm{coh}:\mathcal P^{\mathrm s}\to\mathcal M^{\mathrm s}$ sends $p$
to the class of the cohomology representation of $\mathsf E(p)$. In
particular the form is the same for all holomorphic families of complexes of
the class \eqref{eq:complex} with the same family of cohomology
representations, and it vanishes in every direction along which the
isomorphism class of the cohomology representation is constant.
\end{theorem}

\begin{proof}
$\ker D(p)$ is a holomorphic subbundle. In a local holomorphic frame
$(e_i(p))$ of it let $b(p)\in\mathcal R$ be the matrices of
$A_j(p)|_{\ker D(p)}$; then $b$ is holomorphic and
$\mathrm{coh}=\mathrm{pr}\circ b$. Write
$\iota_{D(p)}(\varpi_V\otimes\varpi_U^{-1})=a(p)\,e_1(p)\wedge\dots\wedge e_{n_0}(p)$
with $a$ holomorphic and nowhere zero. By Theorem~\ref{thm:infimum},
$\cK_\infty(p)=\varkappa(b(p))+\theta_0\log|a(p)|^2+\theta_0\gamma r$, and the
last two terms are pluriharmonic. Apply Lemma~\ref{lem:kappa}(b).
\end{proof}

\begin{remark}\label{rem:levi}
For fixed metrics on $V,U,W$ the function $(A,D)\mapsto\cK(h)$ is not
plurisubharmonic: for the symmetric charges it has a negative Levi direction
along a family of complexes with constant cohomology representation
\cite[Proposition~7.5]{wue26}. Theorem~\ref{thm:moduli} concerns the function
after minimizing over the metrics, where the infimum is not attained; along
such a family it is pluriharmonic.
\end{remark}

\subsection{Convergence along the flow in families}
For a continuous family of initial metrics $h^{\mathrm{in}}(p)$ let $h(t;p)$ be
the trajectory of the phase flow of $\mathsf E(p)$ starting at
$h^{\mathrm{in}}(p)$, and $\cK_t(p)=\cK_{\mathsf E(p)}(h(t;p);\varpi)$.

\begin{theorem}[Convergence in families]\label{thm:families}
Let the charges be ordered.
\begin{enumerate}
\item[(i)] If the cohomology representation of $\mathsf E(p)$ is semistable,
 then $\cK_t(p)\to\cK_\infty(p)$ as $t\to\infty$, although the infimum is not
 attained and no zero of the moment map need exist. For admissible charges
 $\cK_t(p)$ is strictly decreasing in $t$.
\item[(ii)] On compact subsets of $\mathcal P^{\mathrm s}$,
 $0<\cK_t-\cK_\infty\leq Ce^{-ct}$ with $C,c>0$ depending only on the subset.
\item[(iii)] On $\mathcal P^{\mathrm s}$ all parameter derivatives of $\cK_t$
 converge to those of $\cK_\infty$ in the sense of distributions; in
 particular $\ii\partial\bar\partial\cK_t\to\mathrm{coh}^*\omega_\theta$
 weakly.
\end{enumerate}
\end{theorem}

\begin{proof}
(i) By Theorem~\ref{thm:infimum},
$\cK_t-\cK_\infty=\mathcal E(t)+\theta_0\tr{}'E_1(D^\dagger D)+\tr Y$ with
$\mathcal E=\cK_{\coh}-\cK_*\geq0$ as in Section~\ref{sec:convergence}. The
second term is at most $\theta_0rE_1(\lambda(t))$ and the third at most
$s(t)$; both tend to zero by
Theorem~\ref{thm:cancellation}, and $\mathcal E(t)\to0$ by
Theorem~\ref{thm:relaxation}. Monotonicity for admissible charges is
Theorem~\ref{thm:monotone}; positivity of the difference is
Corollary~\ref{cor:infimum}.

(ii) For $t\geq t_1$, $\tr Y\leq s\leq s_1e^{-\frac{\alpha_1}2(t-t_1)}$ and
$\theta_0\tr{}'E_1\leq\theta_0re^{-\lambda}\leq\theta_0re^{\mathfrak c_3}e^{-\frac{\alpha_1}4(t-t_1)}$
by Theorem~\ref{thm:cancellation} and \eqref{eq:E1bounds}, and
$\mathcal E(t)\leq\mathfrak C_2e^{-2\beta t}$ by the proof of
Theorem~\ref{thm:convergence}. All constants there and in
\eqref{eq:constants} are continuous functions of $N(0)$, $\lambda(0)$,
$\tr E_1(L)|_{t=0}$, $\mathcal E(0)$ and $\nu_*$, which are continuous in $p$
for continuous initial metrics; and $\nu_*(p)$ is positive on
$\mathcal P^{\mathrm s}$, being the least eigenvalue of the Hessian form of
Lemma~\ref{lem:gap} at a zero of the moment map, which depends continuously
on $p$ modulo $\mathrm K$ by the proof of Lemma~\ref{lem:kappa}(b), on the
orthogonal complement of the scalars.

(iii) Locally uniform convergence of continuous functions implies convergence
of all derivatives in the sense of distributions.
\end{proof}

\begin{corollary}\label{cor:continuity}
In Theorem~\ref{thm:convergence} the limit $\mathfrak h_\infty$ depends continuously on
the initial metric. Along a continuous family of complexes whose cohomology
representations stay stable, with continuous initial metrics, it depends
continuously on the parameter, in continuous local frames of $\ker D$.
\end{corollary}

\begin{proof}
The constants of Theorem~\ref{thm:convergence} are locally uniform in such
families: $N(0)$, $t_2$ and $\mathcal E(0)$ are continuous, and $\nu_*$ is
continuous and positive on the stable locus, as shown in the proof of
Theorem~\ref{thm:families}(ii). So $d(\mathfrak h(t),\mathfrak h_\infty)\leq\mathfrak Ce^{-\beta t}$ uniformly,
and $\mathfrak h(T)$ depends continuously on the data for each fixed $T$.
\end{proof}


We do not know whether the parameter derivatives of $\cK_t$ converge
pointwise. Outside the stable locus the infimum need not be differentiable.
For example let $k=2$, $n_0=n_1=2$ (so $\theta_0=\theta_1$) and
$b=(\id_2,\left(\begin{smallmatrix}0&a\\a'&0\end{smallmatrix}\right))$ with
$a,a'\neq0$. This representation is isomorphic to the direct sum of the two
representations $(1,\pm c)$ of dimension vector $(1,1)$, $c^2=aa'$, by a
simultaneous conjugation, which does not change $\varkappa$ because
$\theta_0=\theta_1$; and on a direct sum of representations $(1,c_i)$ the
infimum is attained at a diagonal zero of the moment map. Hence
\[
 \varkappa(b)=2\theta_0+2\theta_0\log\frac{1+|a|\,|a'|}{\theta_0}\qquad(a,a'\neq0),
\]
which has no differentiable extension to $a=0$, where the representation is
semistable and not polystable.
